\section{The continuum gap from the single-cube datum}\label{sc:sec:continuum}

We now apply the single-cube ascent to the continuum family of
Definition~\ref{def:family}. The phase hypothesis \textup{(W)} is precisely
that of Definition~\ref{def:WUD}; no particular proposed construction is
used. Its fixed-stencil insertions may have divergent normalisations.
The cube datum controls the thermal trace, so these normalisations do not
enter the partition-function estimate. They enter only separated smears,
whose convergence is already part of \textup{(W)}.

\begin{theorem}[the single-cube continuum criterion]\label{sc:thm:continuum}
Let a continuum family satisfy \textup{(W)}. Suppose there is $t_0>0$ such
that for every sufficiently large $p$ there is $N_0(p)$ such that, for every
$N\ge N_0(p)$ on the family's mesh sequence, the Wilson system at
$\beta_{N,p}$ satisfies \Mcirc{} at an integer $T_0$ with
$3\le T_0\le Nt_0$. Set
\[
\Delta=t_0^{-1}\log\frac{63}{32}>\frac{0.677}{t_0}.
\]
For $L_p\ge2t_0$ and $N\ge N_0(p)$, write
$\mathcal Z_{N,p}=\mathcal Z_{(NL_p)^3}(NL_p)$. Then:
\begin{enumerate}[label=\textup{(\roman*)},leftmargin=*]
\item $\mathcal Z_{N,p}-1\le4096^{-1}e^{-\Delta L_p/2}$;
\item every eigenvalue $\lambda\ne\lambda_0$ of the cylinder transfer matrix
of spatial side $NL_p$ satisfies $0\le\lambda\le e^{-\Delta/N}\lambda_0$;
\item all phase Schwinger functions of \textup{(W)} coincide, and the
Hamiltonian of that Wightman theory satisfies
$\sigma(H)\subset\{0\}\cup[\Delta,\infty)$ and $\ker H=\mathbb C\Omega$.
\end{enumerate}
The theorem does not establish \Mcirc{} along a weak-coupling tuning,
\textup{(W)}, or the nontriviality of the resulting theory.
\end{theorem}

\begin{proof}
Put $n=NL_p$ and $u=32/63$. Since $n\ge2Nt_0\ge2T_0$, purity conversion
and Theorem~\ref{sc:thm:ascent} apply at the actual coupling $\beta_{N,p}$.
They give
\[
\mathcal Z_{N,p}-1\le\frac{u^{n/(2T_0)}}{4096}
\le\frac{u^{L_p/(2t_0)}}{4096}
=\frac{e^{-\Delta L_p/2}}{4096},
\qquad
r_1(n^3)\le u^{1/T_0}\le u^{1/(Nt_0)}=e^{-\Delta/N}.
\]
The strict bound $r_1<1$ also proves simplicity of the top eigenvalue.
This proves (i) and (ii). It remains to identify a physical spectral gap
rather than a statement in lattice units.

\emph{Thermal states and renormalised insertions.}
Let $c$ be a fixed cluster. Translate it in physical time, by an integer
amount independent of $N,p$, so that all its time coordinates are positive;
call the translated cluster $c_0$. At large $N$, its lattice function $G$ is
bounded and lies in a slab $[0,w]$, with $w\le CN$ for a constant depending
only on this fixed cluster and its finite stencils. At large $p$, the
symmetric torus has $M=NL_p\ge2w+2$. Lemma~\ref{lem:thermaldiscrepancy}
gives
\begin{equation}\label{sc:eq:thermal-transfer}
\bigl|S_{N,p}(c)-S^0_{N,p}(c)\bigr|^2
\le(\mathcal Z_{N,p}-1)
\left(S_{N,p}(c_0^\vartheta c_0)-|S_{N,p}(c)|^2\right).
\end{equation}
The bracket is nonnegative. Reflection separates the two sets of insertion
times, so $c_0^\vartheta c_0$ is a legitimate separated cluster. Under
\textup{(W)} the bracket has a finite iterated limit, first in $N$ and then
in $p$. Part (i) tends to zero uniformly in $N$ at large $p$. Thus the
zero-temperature and symmetric-torus smears have the same iterated limit,
in the precise upper-limit sense of Proposition~\ref{prop:ZgivesU}.
This is the weak order-of-limits conclusion \textup{(U$_w$)} used in
Corollary~\ref{cor:Z}; it is not a claim of convergence of the cylinder
smears at each fixed $p$.

\emph{The limiting decay inequality.}
For two clusters $c,d$ whose tests have strictly ordered positive time support, let $F,G$ be their
lattice functions. Their fixed stencils fit in slabs once $N$ is large.
By (ii) and Lemma~\ref{lem:gaptodecay}, for integers
$\sigma\ge\sigma_1\ge1$,
\[
\bigl|S^0_{N,p}(c^\vartheta\tau_\sigma d)
-\overline{S^0_{N,p}(c)}S^0_{N,p}(d)\bigr|
\le e^{-\Delta(\sigma-\sigma_1)}
\sqrt{S^{0,c}_{N,p}(c;\sigma_1)S^{0,c}_{N,p}(d;\sigma_1)},
\]
where $S^{0,c}_{N,p}(c;s)=S^0_{N,p}(c^\vartheta\tau_s c)
-|S^0_{N,p}(c)|^2\ge0$.
Every smear in this inequality is separated. Applying
\eqref{sc:eq:thermal-transfer} to its finitely many clusters, and using
convergence in \textup{(W)}, passes the inequality to their common limit
exactly as in the proof of Corollary~\ref{cor:Z}. With
$\bar S(c)=\int_XS^\chi(c)\,d\nu(\chi)$, it reads
\begin{equation}\label{sc:eq:limit-decay}
\left|\bar S(c^\vartheta\tau_\sigma d)
-\overline{\bar S(c)}\bar S(d)\right|
\le e^{-\Delta(\sigma-\sigma_1)}
\sqrt{\bar S^c(c;\sigma_1)\bar S^c(d;\sigma_1)}.
\end{equation}
For example, the passage does not interchange two uncontrolled limits:
\eqref{sc:eq:thermal-transfer} bounds each replacement error in upper
limit by a constant times $e^{-\Delta L_p/4}$, and that constant remains
bounded as $p\to\infty$ by \textup{(W)}.

\emph{Phases.}
For each such ordered-time cluster $c$ and almost every $\chi$, Lemma~\ref{lem:wightman} identifies
$S^{\chi,c}(c;\sigma)$ with
$\|e^{-\sigma H_\chi/2}\psi'_\chi(c)\|^2\ge0$.
Consequently
\[
\bar S^c(c;\sigma)
=\int_X S^{\chi,c}(c;\sigma)\,d\nu(\chi)
+\operatorname{Var}_\nu(S^\chi(c))
\ge\operatorname{Var}_\nu(S^\chi(c)).
\]
The variance is finite, since its second moment is bounded by the
integrable reflected two-cluster function supplied by \textup{(W)}.
Equation~\eqref{sc:eq:limit-decay}, with $d=c$ and $\sigma_1=1$, makes the
left side tend to zero as $\sigma\to\infty$. Thus $S^\chi(c)$ is almost
surely independent of $\chi$. Use the countability of the tokens and a
countable dense set of tests on each ordered-time configuration space to
choose one null set for all these tests. Smoothness in the configuration
variables, time translation and permutation symmetry then identify the
phase kernels on every separated configuration. Continuity in $\chi$
and full support of $\nu$ extend equality from almost every phase to
every phase, as in Step~3 of Section~\ref{ssec:assembly}.

\emph{Spectrum.}
The common Schwinger functions are those of the Wightman theory in
\textup{(W)}. For a positive-time Euclidean cluster vector put
$\psi'=\psi(c)-\langle\Omega,\psi(c)\rangle\Omega$. By
Lemma~\ref{lem:wightman}, \eqref{sc:eq:limit-decay} gives
\[
\langle\psi',e^{-\sigma H}\psi'\rangle
\le e^{-\Delta(\sigma-1)}\|\psi'\|^2
\quad(\sigma\ge1\text{ an integer}).
\]
If its spectral measure is $\kappa$, then for $0<\eta<\Delta$,
\[
\kappa([0,\Delta-\eta])
\le e^\Delta\|\psi'\|^2e^{-\eta\sigma}\longrightarrow0.
\]
Hence the spectral projection of $[0,\Delta)$ annihilates every such
$\psi'$. These vectors together with $\Omega$ span a dense subspace by
Lemma~\ref{lem:wightman}(iii). The projection is therefore precisely the
projection onto $\mathbb C\Omega$. This proves (iii), including vacuum
uniqueness. The conversion in (ii) is $m_{\rm phys}=Nm_{\rm lat}$, because
$a_N=1/N$; it is this conversion that preserves the positive constant
$\Delta$.
\end{proof}

\subsection{What remains for the Clay statement}\label{sc:ssec:clay}

\begin{figure}[htbp]
\centering
\begin{tikzpicture}[node distance=9mm, every node/.style={align=center,font=\small},
 box/.style={draw,rounded corners,text width=.77\textwidth,inner sep=6pt}]
\node[box] (datum) {Actual Wilson measures along $\beta\to\infty$:\\
$\psi_{(2T_0)^3}(T_0)\le1/4097$, $T_0\le Nt_0$\\
\textbf{Unproved input}};
\node[box,below=of datum] (ascent) {Purity conversion $\longrightarrow$ box comparison $\longrightarrow$ ascent\\
$\mathcal Z_{N,p}-1\le4096^{-1}e^{-\Delta L_p/2}$,
$r_1\le e^{-\Delta/N}$\\\textbf{Proved implication}};
\node[box,below=of ascent] (limit) {A compatible nontrivial Wightman continuum theory\\
\textbf{Existence input, not supplied by the ascent}};
\node[box,below=of limit] (gap) {$\sigma(H)\subset\{0\}\cup[\Delta,\infty)$,
$\ker H=\mathbb C\Omega$, $\Delta=t_0^{-1}\log(63/32)$\\
\textbf{Conditional continuum conclusion}};
\draw[->,thick] (datum)--(ascent);
\draw[->,thick] (ascent)--(limit);
\draw[->,thick] (limit)--(gap);
\end{tikzpicture}
\caption{The exact logical role of the single-cube criterion. The two
boxes marked as inputs are independent obligations. Finite arithmetic
checks and proofs of the arrows do not verify these inputs.}
\label{sc:fig:clay-map}
\end{figure}

The Clay problem asks for a nontrivial theory on $\mathbb R^4$ for every
compact simple gauge group, with the required axioms and a positive mass
gap. The preceding theorem proves an implication for each such group;
it does not prove its antecedents. In particular:
\begin{enumerate}[leftmargin=*]
\item The datum must hold for the full Wilson thermal trace at the
actual weak couplings and with a physical scale bounded independently
of the cutoff. A selected charged sector or a fixed finite lattice
certificate is insufficient.
\item A nontrivial renormalised continuum theory must be constructed and
identified with the scaling limit. Reflection positivity and compactness
of some separated smears do not, by themselves, prove full rotation
invariance, the reconstruction growth bounds, or nontriviality.
\item These two obligations must be discharged for every compact simple
group, including the relevant quotient groups and their flux sectors.
The fact that Proposition~\ref{sc:prop:wilson} covers a group does not
prove the one-cube datum for it.
\end{enumerate}
There is no proof here that the first or second obligation holds along a
four-dimensional weak-coupling tuning. The strengthened criterion and
its verified constants therefore do not constitute an unconditional
solution of the Millennium Prize problem.
